\textbf{Phần 1. Trắc nghiệm lựa chọn}
\setcounterpageref{ex}{0}

\begin{ex} %[3H13N69]
Tháng 8 có bao nhiêu ngày?
\choice
{$30$ ngày}
{\True $31$ ngày}
{$28$ ngày}
{$29$ ngày}
\loigiai{
Tháng 8 có 31 ngày.\\
\textbf{Đáp án của câu hỏi là B}
}
\end{ex}

\begin{ex} %[3H13N69]
Đồng hồ điện tử chỉ $18:15$ tương ứng với
\choice
{$8$ giờ $15$ phút tối}
{$8$ giờ $15$ phút sáng}
{$6$ giờ $15$ phút sáng}
{\True $6$ giờ $15$ phút chiều}
\loigiai{
$18:15$ là $6$ giờ $15$ phút chiều.\\
\textbf{Đáp án của câu hỏi là D}
}
\end{ex}

\begin{ex} %[3H13H69]
Mua 2 gói bánh giá $25,000$ đồng/gói và đưa tờ $100,000$ đồng. Tiền thối lại là
\choice
{$40,000$ đồng}
{$60,000$ đồng}
{\True $50,000$ đồng}
{$70,000$ đồng}
\loigiai{
$100,000 - 25,000 \times 2 = 50,000$ đồng.\\
\textbf{Đáp án của câu hỏi là C}
}
\end{ex}

\begin{ex} %[3H13H69]
Một tiết mục xiếc kéo dài từ 20 giờ đến 20 giờ 45 phút. Tiết mục đó diễn ra trong
\choice
{\True $45$ phút}
{$40$ phút}
{$50$ phút}
{$35$ phút}
\loigiai{
Thời gian: 45 phút.\\
\textbf{Đáp án của câu hỏi là A}
}
\end{ex}

\begin{ex} %[3H13H69]
Nếu ngày 19 tháng 5 là Chủ nhật thì ngày 26 tháng 5 là
\choice
{$Thứ Bảy$}
{\True $Chủ nhật$}
{$Thứ Hai$}
{$Thứ Sáu$}
\loigiai{
$19 + 7 = 26$ nên vẫn là Chủ nhật.\\
\textbf{Đáp án của câu hỏi là B}
}
\end{ex}

\begin{ex} %[3H13H69]
Đổi tờ $50,000$ đồng thành các tờ $10,000$ đồng thì được
\choice
{$4$ tờ}
{$6$ tờ}
{\True $5$ tờ}
{$10$ tờ}
\loigiai{
$50,000 : 10,000 = 5$ tờ.\\
\textbf{Đáp án của câu hỏi là C}
}
\end{ex}

\begin{ex} %[3H13V69]
Mẹ đi chợ mua cá hết $45,000$ đồng, mua rau hết $15,000$ đồng và mua thịt hết $35,000$ đồng. Mẹ đã tiêu hết số tiền là
\choice
{\True $95,000$ đồng}
{$90,000$ đồng}
{$100,000$ đồng}
{$85,000$ đồng}
\loigiai{
$45,000 + 15,000 + 35,000 = 95,000$ đồng.\\
\textbf{Đáp án của câu hỏi là A}
}
\end{ex}

\begin{ex} %[3H13V69]
Một chuyến xe buýt xuất bến lúc 7 giờ 15 phút sáng và đến nơi sau 50 phút. Xe buýt đến nơi lúc mấy giờ?
\choice
{$8$ giờ đúng}
{$8$ giờ $10$ phút}
{$7$ giờ $55$ phút}
{\True $8$ giờ $5$ phút}
\loigiai{
$7$ giờ $15$ phút $+ 50$ phút $= 7$ giờ $65$ phút $= 8$ giờ $5$ phút.\\
\textbf{Đáp án của câu hỏi là D}
}
\end{ex}

\begin{ex} %[3H13V69]
Tháng 11 có bao nhiêu ngày?
\choice
{$31$ ngày}
{\True $30$ ngày}
{$28$ ngày}
{$29$ ngày}
\loigiai{
Tháng 11 có 30 ngày.\\
\textbf{Đáp án của câu hỏi là B}
}
\end{ex}

\begin{ex} %[3H13V69]
Có 3 tờ $20,000$ đồng và 2 tờ $10,000$ đồng. Tổng số tiền là
\choice
{$70,000$ đồng}
{$90,000$ đồng}
{\True $80,000$ đồng}
{$60,000$ đồng}
\loigiai{
$20,000 \times 3 + 10,000 \times 2 = 60,000 + 20,000 = 80,000$ đồng.\\
\textbf{Đáp án của câu hỏi là C}
}
\end{ex}

\begin{ex} %[3H13V69]
Nếu ngày 29 tháng 8 là thứ Sáu thì ngày 2 tháng 9 cùng năm đó là thứ mấy?
\choice
{$Thứ Hai$}
{$Thứ Tư$}
{$Chủ nhật$}
{\True $Thứ Ba$}
\loigiai{
Tháng 8 có 31 ngày. 29 (T6) $\rightarrow$ 30 (T7) $\rightarrow$ 31 (CN) $\rightarrow$ 1/9 (T2) $\rightarrow$ 2/9 (T3).\\
\textbf{Đáp án của câu hỏi là D}
}
\end{ex}

\begin{ex} %[3H13C69]
An có $100,000$ đồng. An mua 2 quyển sách giống nhau và còn thừa $16,000$ đồng. Giá tiền mỗi quyển sách là
\choice
{\True $42,000$ đồng}
{$40,000$ đồng}
{$44,000$ đồng}
{$45,000$ đồng}
\loigiai{
2 quyển sách giá: $100,000 - 16,000 = 84,000$ đồng.\\
Mỗi quyển giá: $84,000 : 2 = 42,000$ đồng.\\
\textbf{Đáp án của câu hỏi là A}
}
\end{ex}

\textbf{Phần 2. Trắc nghiệm đúng sai}
\setcounterpageref{ex}{0}

\begin{ex} %[3H13H69]
Xét tính đúng sai về lịch và thời gian:
\choiceTF
{\True Tháng 7 có 31 ngày.}
{\True Tháng 8 có 31 ngày.}
{Tháng 9 có 31 ngày.}
{\True Tháng 9 có 30 ngày.}
\loigiai{
\begin{itemize}
    \item Tháng 7 có 31 ngày: Đúng.
    \item Tháng 8 có 31 ngày: Đúng.
    \item Tháng 9 có 30 ngày (không phải 31). Ý 3 Sai, Ý 4 Đúng.
\end{itemize}
\begin{itemize}
    \item a) \textbf{Đúng}.
    \item b) \textbf{Đúng}.
    \item c) \textbf{Sai}.
    \item d) \textbf{Đúng}.
\end{itemize}
}
\end{ex}

\begin{ex} %[3H13V69]
Mẹ có tờ $100,000$ đồng. Mua $1\text{ kg}$ cam giá $35,000$ đồng và $1\text{ kg}$ nho giá $45,000$ đồng. Xét tính đúng sai:
\choiceTF
{\True Tổng số tiền mua hoa quả là $80,000$ đồng.}
{\True Số tiền mẹ còn lại là $20,000$ đồng.}
{Tiền mua nho nhiều hơn tiền mua cam là $5,000$ đồng.}
{\True Tiền mua nho nhiều hơn tiền mua cam là $10,000$ đồng ($45000 - 35000 = 10000$).}
\loigiai{
\begin{itemize}
    \item Tổng: $35000 + 45000 = 80000$ đồng. Đúng.
    \item Còn lại: $100000 - 80000 = 20000$ đồng. Đúng.
    \item Nho hơn cam: $10000$ đồng $\neq 5000$. Ý 3 Sai, Ý 4 Đúng.
\end{itemize}
\begin{itemize}
    \item a) \textbf{Đúng}.
    \item b) \textbf{Đúng}.
    \item c) \textbf{Sai}.
    \item d) \textbf{Đúng}.
\end{itemize}
}
\end{ex}

\textbf{Phần 3. Trả lời ngắn}
\setcounterpageref{ex}{0}

\begin{ex} %[3H13H69]
Tháng 10 có bao nhiêu ngày?
\shortans{31}
\loigiai{
\textbf{Đáp số:} $31$\\
31 ngày.
}
\end{ex}

\begin{ex} %[3H13H69]
Một giờ có bao nhiêu phút?
\shortans{60}
\loigiai{
\textbf{Đáp số:} $60$\\
60 phút.
}
\end{ex}

\begin{ex} %[3H13V69]
Mua 4 gói kẹo giá $15,000$ đồng/gói. Tính tổng số tiền phải trả theo đồng.
\shortans{60000}
\loigiai{
\textbf{Đáp số:} $60000$\\
$15000 \times 4 = 60000$ đồng.
}
\end{ex}

\begin{ex} %[3H13V69]
Nếu ngày 10 tháng 10 là thứ Năm thì ngày 24 tháng 10 cùng năm là thứ mấy (ghi số thứ, ví dụ thứ Năm ghi 5)?
\shortans{5}
\loigiai{
\textbf{Đáp số:} $5$\\
$10 + 14 = 24$ (sau đúng 2 tuần) nên vẫn là thứ Năm (ghi 5).
}
\end{ex}

\begin{ex} %[3H13V69]
Một buổi chiếu phim bắt đầu lúc 19 giờ 15 phút và kết thúc lúc 21 giờ. Buổi chiếu phim kéo dài bao nhiêu phút?
\shortans{105}
\loigiai{
\textbf{Đáp số:} $105$\\
1 giờ 45 phút $= 105$ phút.
}
\end{ex}

\begin{ex} %[3H13C69]
Bạn Bình có $100,000$ đồng. Bình mua 3 quyển vở cùng loại và được thối lại $25,000$ đồng. Hỏi mỗi quyển vở giá bao nhiêu đồng?
\shortans{25000}
\loigiai{
\textbf{Đáp số:} $25000$\\
3 quyển vở giá: $100000 - 25000 = 75000$ đồng.\\
Mỗi quyển: $75000 : 3 = 25000$ đồng.
}
\end{ex}
